% 本文件由 LaTeX 排版 Agent 根据有效 translation.json 自动生成。
% author=Tertullian; work_id=0319; title=驳斥一切异端

\chapter{驳斥一切异端}

% structure: p0001 source=h1 target=omitted reason=page-title-already-rendered-by-parent

% structure: p0002 source=h2 target=section reason=relative-heading-rank; base=h2

\section{第一章 最早的异端者：\Person{西满·玛古斯}、\Person{默纳德尔}、\Person{撒督尼努}、\Person{巴西理德}、\Person{尼各老}。〖作品以残篇开始〗}

关于这些异端者，我将（略过许多）简要归纳几点。因为对于犹太教的异端者，我保持沉默——我是指撒玛黎雅人多西特乌斯，他是第一个敢于否认先知的人，理由是他们并非在圣神的默感下说话。对于撒杜塞人，我保持沉默，他们从这一错误的根源出发，进而胆敢否认肉身的复活。法利塞人，我略去不谈，他们因在法律上附加某些增补而与犹太人 \Quote{分岐}，这也使他们配得上这一名称；与他们一起的，还有黑落德党人，他们声称黑落德就是基督。我转向那些选择以福音作为他们异端起点的异端者。\par

其中第一个便是\Person{西满·玛古斯}，他在\BibleBook{宗徒}{大事录}中受到宗徒伯多禄公正而应得的判决。他竟敢称自己为至高德能，即至高天主；并（声称）宇宙由其天使所造；他降来寻找一个游荡的魔鬼，即智慧；他以天主的幻像，在犹太人中间 \Emph{没有} 受难，而是 \Emph{仿佛受了难}。\par

在他之后，其门徒\Person{默纳德尔}（同样是一个术士）说与西满相同的话。凡西满自称的，默纳德尔也同样自称，断言除非因他的名受洗，否则无人能得救恩。\par

随后，\Person{撒督尼努} 接踵而至；他也声称非受生的德能，即天主，居于最高区域，那些区域是无限的，且在紧邻我们之上的区域；但远离祂的天使制造了下界；并且因为从上面来的光在下界闪耀照耀，天使们便小心地试图按那光的肖像造人；那人爬行在地面上；这光与这更高的德能，由于慈悲，成为人内可蒙救恩的火花，而其余部分尽皆消灭；基督并无实体存在，仅以幻形忍受一种\OriginalTerm{准受难}{quasi-passion}；肉身的复活绝不会发生。\par

此后，异端者\Person{巴西理德} 爆发。他声称有一位至高神祇，名为阿布拉克萨斯，由祂创造了心智，希腊文称为\OriginalTerm{努斯}{\GreekText{Νοῦς}}；从那里生出圣言；由圣言生出眷顾、德能与智慧；从这些随后生出宰制者、掌权者与天使；由此产生了无穷的天使流出与涌现；由这些天使制造了三百六十五重天与世界，为荣耀阿布拉克萨斯，其名若计算，本身便含有此数。现在，在这些制造世界的最后天使中，他将犹太人的天主列为最后，即法律与先知的天主，他否认祂是天主，却声称是一位天使。他说，亚巴郎的后裔分配给了祂，因此祂带领以色列子民从埃及地进入客纳罕地；他声称祂比其他天使更好乱，因此常常激起叛乱与战争，甚至流人的血。此外，他声称基督并非由这位世界的创造者所派遣，而是由上述的阿布拉克萨斯派遣；祂以幻影而来，毫无肉体实质；不是祂在犹太人中间受难，而是西满被钉死以代替祂；因此，不可相信那位被钉者，免得人承认相信了西满。他说，殉道不应忍受。他极力攻击肉身的复活，声称救恩并未应许给 \Emph{身体}。\par

又一位异端兄弟是\Person{尼各老}。他是\BibleBook{宗徒}{大事录}中被选立的七位执事之一。他声称黑暗对光明起了贪欲——而且是污秽丑陋的贪欲；从这种混杂中，羞于提及的腐臭不洁之物（产生出来）。其余（教义）也是污秽的。因为他讲述某些永世，即淫乱之子，以及可憎可耻的拥抱与混杂的结合，还有某些更加卑劣的产物。他还教导生出了魔鬼、诸神、七灵，以及其他足够亵渎且污秽的事物，我们羞于叙述，立即略过。对于我们而言，这一尼苛劳党人的异端已被主的\BibleBook{}{默示录}以最权威的判决所谴责，其中说：\BibleQuote{你之所以持守这个，因为你恨恶尼苛劳党人的教训，我也恨恶。}\par

% structure: p0009 source=h2 target=section reason=relative-heading-rank; base=h2

\section{第二章 蛇派、加音派、舍特派}

另外，还有那些被称作\Emph{蛇派}的异端者也加入其中：因为他们如此推崇蛇，甚至认为蛇比基督本人更优越；因为他们说，正是蛇给了我们善恶知识的起源。梅瑟领悟到蛇的权能与威严（他们说），便铸造了铜蛇；凡注视它的人都得到了医治。基督本人（他们进一步说）在祂的福音中效法了梅瑟之蛇的神圣权能，说道：\BibleQuote{就如梅瑟在旷野里举起了蛇，人子也必须照样被举起来。}\BibleRef{若 3:14}他们引进蛇来祝福他们的圣体元素。如今，这一谬误的全部排场和教训源自以下来源。他们说，从至高的首要\OriginalTerm{永世}{Æon}（\Emph{人们所说的}）中流溢出几个较低等的永世。然而，有一位名叫\Emph{雅达巴沃特}的永世对抗所有这些永世。他是由第二位永世与较低等永世混合而怀有的；后来，当他渴望强行进入更高领域时，却因自身与物质的沉重混合而无法到达更高领域；他被留在中间，并伸展到他的全部大小，从而制造了苍穹。但是，\Emph{雅达巴沃特}下降到更低处，为自己生了七个儿子，并通过自我扩展封闭了他们向上看的视线，以便（这些）天使无法知道上面是什么，从而可能认为他是唯一的天主。因此，这些低等的权能和天使制造了\Emph{人}；由于人是由较弱和中等的力量起源的，他像虫子一样爬行。然而，雅达巴沃特所出自的那个永世，心中嫉妒，向爬行的人注入了一丝火花；人受此激发，将因谨慎而变得智慧，并能够理解上面的事物。同样，前述的\Emph{雅达巴沃特}愤怒地，从他自身发出了蛇的权能及相似形；这（权能）就是在乐园中的权能——也就是说，这就是\Emph{蛇}——厄娃相信了它，仿佛它是天主圣子。他们说，它从树上摘了果子，并将善恶知识赐给了人类。此外，基督并非存在于血肉的实体中：肉身的救恩是根本无望的。\par

此外，还有另一个异端也出现了，即所谓的\Emph{加音派}。其原因是，他们推崇加音，仿佛他是被某种运作在他内的强大权能所怀有的；因为亚伯尔是由一个较低等的权能怀有而出生，因此被证明是低等的。主张此说者也为叛徒犹达斯辩护，告诉我们他是可敬而伟大的，因为他所夸耀的施予人类的益处；因为他们中有些人认为，应当为此感谢犹达斯：也就是说，他们说，犹达斯观察到基督想要颠覆真理，便出卖了祂，以便真理不可能被颠覆。而另一些人如此反驳他们，说：因为今世的权势们不愿意基督受难，唯恐祂的死亡为人类预备救恩，他（犹达斯）为了人类救恩着想，出卖了基督，以便救恩绝无可能被阻碍——这救恩当时正被那些反对基督受难的权能所阻碍；这样，通过基督的受难，人类救恩就绝无可能被延迟。\par

但是，又有异端兴起，即所谓的\Emph{舍特派}。该谬误的教义如下。两位人类由天使所造——加音和亚伯尔。因他们，天使中发生了巨大的争斗与不和；为此，那超越所有权能的权能——他们称之为\OriginalTerm{母亲}{Mother}——当他们说亚伯尔被杀时，愿意他们的这位舍特被怀有和出生以代替亚伯尔，以便那些创造了先前那两位人类的天使被剥夺，而此纯净的种子兴起并出生。因为，他们说，曾有两位天使与人类不义的混合；因此，那（如我们所说）他们称之为母亲的权能甚至带来洪水，为报仇，以便那混合的种子被扫除，而此唯一纯净的种子被完整保存。但是（徒劳）：因为那些产生先前种子的人，秘密而隐秘地（且不为那母亲权能所知）送入方舟内，连同那些\BibleQuote{八条灵魂}以及含的种子，以便邪恶的种子不致灭亡，反而与其余一起被保存，并在洪水后恢复到地上，并效法其余，生长、扩散、充满并占据整个大地。此外，关于基督，他们的想法如此：他们仅仅称祂为舍特，说祂是代替实际的舍特而来。\par

% structure: p0013 source=h2 target=section reason=relative-heading-rank; base=h2

\section{第三章 卡尔波克拉特、则林托、厄彼雍}

此外，卡尔波克拉特引入了以下派别。他断言，有一位至高权能，是上位（区域）的首要者：从这权能产生了天使和权能，它们远离上位权能，创造了这低下的世界：基督并非由\Emph{童贞}玛利亚所生，而是由若瑟的种子所生——纯然一个人——在正义的实践与生活的完整上超越所有其他人（他们承认）；祂在犹太人中间受难；只有祂的灵魂被接到天上，因为它比所有其他灵魂更坚固更刚毅：由此他推论，只保留灵魂的救恩，认为没有身体的复活。\par

在他之后，异端者则林托兴起，教导类似。因为他也说，世界是由那些\Emph{天使}所起源；并提出基督是由若瑟的种子所生，主张祂纯粹是人，没有天主性；还肯定法律是由天使所赐；认为犹太人的天主不是主，而是一位天使。\par

\Person{厄彼翁}继他而起，并不完全与\Person{色林托}一致；他宣称世界由天主创造，而非天使；并且因为经上记载：\BibleQuote{没有门徒高过\Emph{他的}先生，也没有仆人高过\Emph{他的}主人}，他同样提出律法是\Emph{有约束力的}，当然是为了排除福音并辩护犹太教。\par

% structure: p0017 source=h2 target=section reason=relative-heading-rank; base=h2

\section{第四章 \Person{瓦伦廷}、\Person{托勒密}与\Person{塞昆杜斯}、\Person{赫拉克利翁}}

异端者\Person{瓦伦廷}还引入了许多神话。我将精简这些神话，并作简要概述。因为他引入了\OriginalTerm{圆满}{Pleroma}与三十个\OriginalTerm{永世}{Æons}。此外，他以\OriginalTerm{配对}{syzygies}，即某种婚姻结合的方式解释这些永世。因为在最初的（永世）中，他说，有\OriginalTerm{深渊}{Bythus}与\OriginalTerm{沉寂}{Silence}；从它们生出\Person{理智}与\Person{真理}；从理智与真理涌出\Person{圣言}与\Person{生命}；从圣言与生命又创造了\Person{人}与\Person{教会}。但（这些并非全部）；因为从最后这些也生出了十二个永世；此外，从\Person{言语}与\Person{生命}又\Emph{发出}其他十个永世：这就是由八联、十联与十二联在圆满中构成的三十之数。第三十个永世想要看见伟大的深渊；为了看见他，竟胆敢升到高处；但因不能看见他的伟大而沮丧，几乎要消散，要不是某位——他称之为\OriginalTerm{界限}{Horos}——被派来增强他，通过说出\Quote{亚欧}这个词而坚固了他。这个陷入沮丧的永世，他称之为\OriginalTerm{阿卡莫特}{Achamoth}，并说它被某些悔恨的情感所攫住，并从它的情感中产生了物质性的本质。因为它惊恐、恐惧、并被悲伤所压倒；他从这些情感中受孕并生产。因此他创造了天、地、海以及其中的一切；因此他所造的一切都是软弱、脆弱、易堕落和必死的，因为他自己是从沮丧中受孕并产生的。然而，他从\Person{阿卡莫特}因恐惧、惊恐、悲伤或汗水所提供的那些物质本质中创造了这个世界。因为他宣称，从它的惊恐中产生了黑暗；从它的恐惧与无知中产生了邪恶与恶意的灵；从它的悲伤与眼泪中产生了泉源的水流、江河与海的物质本质。此外，\Person{基督}被那位第一父——即深渊——所派遣。祂并不在我们肉身的实体中；而是从天带来某种属灵的身体，像水通过管子一样经过\Person{童贞玛利亚}，既未从中接受也未从中借用任何东西。祂否认我们现今肉身的复活，但（主张）某种姐妹肉身的复活。对于法律和先知，他赞同一些，不赞同一些；即他不赞同所有，因他否定了一些。他还有一部自己的福音，不同于我们的这些。\par

在他之后兴起了异端者\Person{托勒密}与\Person{塞昆杜斯}，他们完全赞同\Person{瓦伦廷}，仅在以下一点不同：即\Person{瓦伦廷}曾构想三十个永世，他们则增加了几个；他们先加了四个，后又加了四个。而\Person{瓦伦廷}的断言——即第三十个永世离开了圆满（如同陷入沮丧）——他们予以否认；他们说，那个因渴望看见第一父未果而沮丧的永世并不属于原初的三十之数。\par

此外，还兴起了\Person{赫拉克利翁}，一位同道异端者，其观点与\Person{瓦伦廷}配对；但凭借术语上的一些新奇，他想要显得在观点上有所区别。因为他引入了这样的概念：首先存在他所谓的（一个\OriginalTerm{单子}{Monad}）；然后从那个单子（产生了）两个，然后是其余永世。然后他引入了\Person{瓦伦廷}的整个体系。\par

% structure: p0021 source=h2 target=section reason=relative-heading-rank; base=h2

\section{第五章 \Person{玛尔谷}与\Person{科拉巴苏斯}}

在这些之后，不乏一位\Person{玛尔谷}和一位\Person{科拉巴苏斯}，他们用希腊字母组成了一种新颖的异端。因为他们断言，没有这些字母就无法找到真理；不仅如此，在这些字母中包含了真理的全部圆满与完善；因此基督曾说：\BibleQuote{我是\OriginalTerm{阿尔法}{\GreekText{Α}}与\OriginalTerm{欧米伽}{\GreekText{Ω}}}。实际上，他们说耶稣基督降临了，即那\OriginalTerm{鸽子}{\GreekText{περιστερά}}降在了耶稣身上；既然鸽子在希腊文中被称为\OriginalTerm{鸽子}{\GreekText{περιστερά}}，它本身就包含数字八〇一。这些人遍历\OriginalTerm{\GreekText{Ω}}{\GreekText{Ω}}、\OriginalTerm{\GreekText{Ψ}}{\GreekText{Ψ}}、\OriginalTerm{\GreekText{Χ}}{\GreekText{Χ}}、\OriginalTerm{\GreekText{Φ}}{\GreekText{Φ}}、\OriginalTerm{\GreekText{Υ}}{\GreekText{Υ}}、\OriginalTerm{\GreekText{Τ}}{\GreekText{Τ}}——实则遍历整个字母表，直至\OriginalTerm{\GreekText{Α}}{\GreekText{Α}}与\OriginalTerm{\GreekText{Β}}{\GreekText{Β}}——并计算八联与十联。因此我们可以承认，述说他们所有的琐事是无用且无聊的。然而，不仅徒然而且危险的是这一点：他们捏造了第二位神，在造物主之外；他们肯定基督不在肉身的实体中；他们说没有肉身的复活。\par

% structure: p0023 source=h2 target=section reason=relative-heading-rank; base=h2

\section{第六章 \Person{克尔多}、\Person{马尔西翁}、\Person{卢卡努斯}、\Person{阿佩莱斯}}

对此又加上一位\Person{克尔多}。他引入了两个第一原因，即两位神——一位良善，一位残酷；良善者为更高者，后者即残酷者是世界的造物主。他拒绝预言和律法；否认天主造物主；主张前来的基督是那位更高之神之子；肯定祂不在肉身的实体中；说祂仅有幻影形状，并未真正受苦，而是经历了某种准受难，并非由童贞女所生，\Emph{其实}根本未曾降生。他仅赞同灵魂的复活，否认身体的复活。他只接受\WorkTitle{路加福音}，且并非全本。对于宗徒\Person{保禄}的书信，他既不取全部，也不取完整本。他将\WorkTitle{宗徒大事录}和\WorkTitle{默示录}斥为伪书。\par

在他之后出现了一位他的门徒，名为\Person{马尔西翁}，本都人，一位主教的儿子，因对某贞女犯下强暴而被逐出教会。他根据经上所说：\BibleQuote{好树结好果子，坏树结坏果子}，试图赞同\Person{克尔多}的异端；因此他的主张与在他之前的异端者完全相同。\par

在他之后兴起了一位名为\Person{卢卡努斯}的人，是\Person{马尔西翁}的追随者和门徒。他也涉足相同种类的亵渎，教导与\Person{马尔西翁}和\Person{克尔多}所教导的相同。\par

紧跟他们之后的是马西昂的门徒阿培勒；他因陷于自身的肉欲而被马西昂断绝关系。他引进了一位在无限高天之上的天主，并说祂创造了诸多权能与天使；此外，还有另一位德能，他声称被称为\Quote{主}，但将之描述为一位天使。他借此宣称世界是按照一个更高世界的样式而产生的。他将悔改的原则贯穿于这个\Emph{低级}世界中，因为他没有像那个更高世界被创造那样完美地创造它。他弃绝律法和先知。关于基督，他既不像马西昂那样肯定祂是幻影形状，也不像福音所教导的那样具有真实身体的实质，而是说，因为祂从高天降下，在降临过程中为自己编织了星质和空气般的肉体；并在复活时，在上升过程中将祂降临时所借用的各部分归还给了各个元素：这样——祂身体的各部分消散后——祂只将祂的神灵重新安置在天上。此人否认肉身的复活。他也只使用一位宗徒，但那是马西昂的，即残缺的。他只教导灵魂的得救。此外，他还有自己的私密而特别的读物，他称之为一位名叫菲鲁美纳的少女的\Quote{显现}，他追随她为先知。此外，他还有自己的书籍，他称之为\WorkTitle{三段论}的书，在其中他试图证明梅瑟关于天主所写的一切都不是真的，而是假的。\par

% structure: p0028 source=h2 target=section reason=relative-heading-rank; base=h2

\section{第七章 塔提安、卡塔弗里吉亚派、卡塔普罗克兰派、卡塔埃斯基内坦派}

这些异端者中又加上了塔提安，他是异端兄弟。此人是殉道者犹斯定的门徒。犹斯定死后，他开始抱有与他不同的见解。因为他完全带有瓦伦廷的气味；并加上一点：亚当甚至不能获得救恩；仿佛当树枝可得救时，树根却不行！\par

其他异端者扩大了名单，他们被称为卡塔弗里吉亚派，但他们的教导并不统一。因为其中有\Emph{有些}被称为卡塔普罗克兰派；另一些被称为卡塔埃斯基内坦派。他们有一种共同的亵渎，还有一种并非共同而是特殊和特别的亵渎。共同的亵渎在于他们说圣神确实在宗徒们内，但护慰者却不在；以及他们说护慰者在孟他努斯里说了比基督带入福音范围更多的内容，而且不仅更多，而且更好、更大。但跟随埃斯基内斯的人有特别的亵渎，即他们添加了这一条：他们肯定基督本人就是圣子和圣父。\par

% structure: p0031 source=h2 target=section reason=relative-heading-rank; base=h2

\section{第八章 布拉斯特斯、两位狄奥多图斯、普拉克西亚斯}

除了所有这些，还有布拉斯特斯，他企图暗中引入犹太教。因为他说逾越节不可按其他方式，而只能按梅瑟律法在正月十四日守。但谁不会看出，若他将基督带回律法，福音的恩宠就被剥夺了呢？\par

再加上拜占庭人狄奥多图斯，他曾因基督之名被逮捕并背教，此后不停地亵渎基督。因为他引入了一个学说，肯定基督仅仅是一个人，而否认祂的神性；教导说祂确实由圣神从童贞女所生，但只是一个孤独而赤裸的人，在其余人中并无优越之处，只有义德的优越。\par

在他之后，又爆发了第二个异端的狄奥多图斯，他本人又引入了一个姊妹宗派，并说那人基督本人同样是由圣神和童贞玛利亚感孕而生，但祂比默基瑟德低劣；因为关于基督记载说：\BibleQuote{你照默基瑟德的品位，永为司祭。}因为他说，那默基瑟德是一个具有卓越恩宠的天上德能；基督为人类行事，成为他们的恳求者和中保；默基瑟德则为天上的天使和德能行事。因为他说，他优于基督到如此程度，以至于他是\OriginalTerm{无父者}{\GreekText{ἀπάτωρ}}、\OriginalTerm{无母者}{\GreekText{ἀμήτωρ}}、\OriginalTerm{无谱系者}{\GreekText{ἀγενεαλόγητον}}，他的开始和结束都未被领悟，也不能被领悟。\par

但在所有这些之后，又有一个普拉克西亚斯引入了一种异端，维克托里努斯谨慎地加以证实。他断言耶稣基督就是全能的天主圣父。他主张祂被钉十字架、受难并死亡；此外，以亵渎和渎神的鲁莽，他坚持说祂本人正坐在自己的右边。\par

\SourceRecord{Translated by S. Thelwall.；Ante-Nicene Fathers；Vol. 3.；Edited by Alexander Roberts, James Donaldson, and A. Cleveland Coxe.；Buffalo, NY: Christian Literature Publishing Co.,；1885.；<http://www.newadvent.org/fathers/0319.htm>.}
